% Videos dentro del PDF (sin audio): lámina 1 con \movie (paquete multimedia, archivo externo video.mp4), lámina 2 con
% media9 (video incrustado en una anotación RichMedia), lámina 3 con \movie a un archivo que no existe.
\pdfminorversion=4
\pdfobjcompresslevel=0
\documentclass[aspectratio=169]{beamer}
\usepackage{multimedia}
\usepackage{media9}
\setbeamertemplate{navigation symbols}{}
\begin{document}
\begin{frame}{Video externo (\texttt{\textbackslash movie})}
  \centering
  \movie[autostart,loop,width=8cm,height=4.5cm,poster]{\fbox{\rule{0pt}{4.5cm}\rule{8cm}{0pt}}}{video.mp4}
\end{frame}
\begin{frame}{Video incrustado (media9)}
  \centering
  \includemedia[width=8cm,height=4.5cm,activate=pageopen,addresource=video.mp4,flashvars={source=video.mp4&autoPlay=true&loop=true}]{\fbox{\rule{0pt}{4.5cm}\rule{8cm}{0pt}}}{VPlayer.swf}
\end{frame}
\begin{frame}{Video externo que no existe}
  \centering
  \movie[autostart,width=8cm,height=4.5cm]{\fbox{\rule{0pt}{4.5cm}\rule{8cm}{0pt}}}{no_existe.mp4}
\end{frame}
\end{document}
